\chapter{Hito 3: Diseño de las Fuentes de Datos e Ingesta Inicial}

\section{Objetivo}

El objetivo de este hito es construir la primera capa del flujo de datos.

Al finalizar este hito el estudiante será capaz de:

\begin{itemize}
\item Comprender las fuentes de información del proyecto.
\item Diseñar datasets alineados con las necesidades de negocio.
\item Organizar correctamente las carpetas de datos.
\item Crear archivos CSV simulados.
\item Leer datos desde Python.
\item Realizar validaciones básicas antes de cargar a PostgreSQL.
\end{itemize}

\section{¿Por qué comenzamos con archivos CSV?}

En proyectos reales los datos suelen provenir de:

\begin{itemize}
\item APIs.
\item Bases de datos.
\item Excel.
\item CSV.
\item Servicios web.
\item Plataformas regulatorias.
\end{itemize}

Sin embargo, durante las primeras etapas del proyecto es recomendable trabajar con archivos CSV.

Esto permite:

\begin{itemize}
\item Diseñar el modelo de datos.
\item Construir el ETL.
\item Validar la arquitectura.
\item Evitar complejidad innecesaria.
\end{itemize}

Una vez validado el proceso, sustituiremos los CSV por fuentes reales.

\section{Arquitectura de la ingesta}

El flujo de trabajo será:

\begin{verbatim}
Fuente CSV
|
V
data/raw
|
V
Python
|
V
Validación
|
V
data/processed
|
V
PostgreSQL
\end{verbatim}

\section{Fuentes utilizadas}

El proyecto utilizará dos fuentes principales.

\subsection{Fuente 1: Tipo de Cambio}

Representa información similar a la publicada por el Banco Central.

Grano:

\begin{quote}
Una observación por fecha y moneda.
\end{quote}

Variables:

\begin{itemize}
\item fecha
\item moneda
\item tasa\_compra
\item tasa\_venta
\end{itemize}

Ejemplo:

\begin{verbatim}
fecha,moneda,tasa_compra,tasa_venta
2024-01-01,USD,58.50,58.80
2024-01-01,EUR,63.10,63.60
\end{verbatim}

\subsection{Fuente 2: Balances Bancarios}

Representa información similar a la publicada por la Superintendencia de Bancos.

Grano:

\begin{quote}
Una observación por mes y entidad.
\end{quote}

Variables:

\begin{itemize}
\item periodo
\item entidad\_cod
\item entidad\_nombre
\item tipo\_entidad
\item activos\_totales
\item cartera\_creditos
\item depositos
\item patrimonio
\item resultado\_neto
\item depositos\_me
\item cartera\_me
\end{itemize}

Ejemplo:

\begin{verbatim}
periodo,entidad_cod,activos_totales
202401,B001,120000000
\end{verbatim}

\section{Paso 1: Crear estructura de almacenamiento}

Verificar que existe:

\begin{verbatim}
data/

+-- raw/
+-- processed/
\end{verbatim}

\texttt{raw} contendrá los datos originales.

\texttt{processed} contendrá los datos validados.

\section{Paso 2: Crear archivo tipo\_cambio.csv}

Crear:

\begin{verbatim}
data/raw/tipo_cambio.csv
\end{verbatim}

Contenido mínimo:

\begin{verbatim}
fecha,moneda,tasa_compra,tasa_venta
2024-01-01,USD,58.50,58.80
2024-01-02,USD,58.55,58.85
2024-01-03,USD,58.60,58.90
2024-01-01,EUR,63.10,63.60
2024-01-02,EUR,63.20,63.70
\end{verbatim}

\section{Paso 3: Crear archivo balances.csv}

Crear:

\begin{verbatim}
data/raw/balances.csv
\end{verbatim}

Contenido mínimo:

\begin{verbatim}
periodo,entidad_cod,entidad_nombre,
tipo_entidad,activos_totales,
cartera_creditos,depositos,
patrimonio,resultado_neto,
depositos_me,cartera_me

202401,B001,Banco Alfa,
Banco Multiple,120000000,
85000000,95000000,
15000000,1200000,
25000000,18000000
\end{verbatim}

\section{Paso 4: Crear script de exploración}

Crear:

\begin{verbatim}
etl/explore_sources.py
\end{verbatim}

\begin{lstlisting}[language=Python]
import pandas as pd

fx = pd.read_csv(
"data/raw/tipo_cambio.csv"
)

balances = pd.read_csv(
"data/raw/balances.csv"
)

print("\nTipo de Cambio")
print(fx.head())

print("\nBalances")
print(balances.head())
\end{lstlisting}

\section{Paso 5: Ejecutar exploración}

Ejecutar:

\begin{lstlisting}[language=bash]
python etl/explore_sources.py
\end{lstlisting}

Resultado esperado:

\begin{itemize}
\item Visualización de las primeras filas.
\item Confirmación de lectura correcta.
\end{itemize}

\section{Paso 6: Verificar estructura}

Agregar:

\begin{lstlisting}[language=Python]
print(fx.info())

print(balances.info())
\end{lstlisting}

Analizar:

\begin{itemize}
\item Número de filas.
\item Número de columnas.
\item Tipos de datos.
\item Valores faltantes.
\end{itemize}

\section{Paso 7: Validaciones básicas}

Agregar:

\begin{lstlisting}[language=Python]
print(
fx.isnull().sum()
)

print(
balances.isnull().sum()
)
\end{lstlisting}

Verificar:

\begin{itemize}
\item Valores nulos.
\item Fechas inválidas.
\item Tasas negativas.
\end{itemize}

\section{Paso 8: Crear capa processed}

Convertir tipos de datos:

\begin{lstlisting}[language=Python]
fx["fecha"] = pd.to_datetime(
fx["fecha"]
)

fx["tasa_compra"] = (
fx["tasa_compra"]
.astype(float)
)

fx["tasa_venta"] = (
fx["tasa_venta"]
.astype(float)
)
\end{lstlisting}

\section{Paso 9: Guardar datos procesados}

\begin{lstlisting}[language=Python]
fx.to_csv(
"data/processed/tipo_cambio.csv",
index=False
)

balances.to_csv(
"data/processed/balances.csv",
index=False
)
\end{lstlisting}

\section{Paso 10: Documentar las fuentes}

Crear:

\begin{verbatim}
docs/data_dictionary.md
\end{verbatim}

Documentar:

\begin{itemize}
\item Nombre de variable.
\item Descripción.
\item Tipo de dato.
\item Fuente.
\item Frecuencia.
\end{itemize}

\section{Entregables}

Al finalizar este hito deberá existir:

\begin{itemize}
\item Archivos CSV en raw.
\item Archivos validados en processed.
\item Script de exploración.
\item Diccionario de datos.
\item Validaciones básicas implementadas.
\end{itemize}

\section{Checklist de validación}

\begin{itemize}
\item[$\Box$] Existe carpeta raw.
\item[$\Box$] Existe carpeta processed.
\item[$\Box$] tipo\_cambio.csv creado.
\item[$\Box$] balances.csv creado.
\item[$\Box$] Python puede leer ambos archivos.
\item[$\Box$] Validaciones ejecutadas.
\item[$\Box$] Archivos procesados generados.
\end{itemize}

\section{Próximo hito}

En el Hito 4 diseñaremos el modelo dimensional.

Construiremos:

\begin{itemize}
\item Dimensión Tiempo.
\item Dimensión Moneda.
\item Dimensión Entidad Financiera.
\item Hecho Tipo de Cambio.
\item Hecho Balance Bancario.
\item Primer esquema estrella.
\end{itemize}
